\section{The cube condition at strong coupling}\label{sec:bpstrong}

Keep the setting of Sec.~\ref{sec:cube}, with a bi-invariant metric $\delta$
of diameter $D$ on $G$, invariant under inversion, and $d_E(U,U')=\sum_e
\delta(U_e,U'_e)$. Two distinct links are \emph{partners} if they share a
plaquette; two links share at most one plaquette, and every link has eighteen
partners, three in each of its six plaquettes. The direction field $\sigma$ of
Sec.~\ref{sec:linefield} has period two, so it is defined on $\Lambda_L$. By
Lemma~\ref{lem:structure}(i) the links of $F$ with an end point $v$ are exactly
$(v,\sigma(v))$ and $(v-e_{\sigma(v)},\sigma(v))$, and $\sigma(v-e_{\sigma(v)})=
\sigma(v)$.

\begin{proposition}\label{prop:bpsup}
At $\beta=0$, $\BP(n,1)$ holds for every $n$.
\end{proposition}

\begin{proof}
For $\varphi\equiv1$, $\mu$ is a product of Haar measures. Write $f=\sum_Sf_S$
with $f_S$ a function of the links of $S$ and $\int f_S\,dU_e=0$ for $e\in S$;
the $f_S$ are orthogonal and $\mathcal P_xf=\sum_{S\cap V_x=\emptyset}f_S$, so
$\sum_x\mu((f-\mathcal P_xf)^2)=\sum_{S\ne\emptyset}\|f_S\|^2\#\{x:S\cap V_x\ne
\emptyset\}\ge n^4\Var_\mu(f)$.
\end{proof}

\subsection*{Gauge fixing inside a cube}
For a cube $Q_x=x+[0,n)^4$ write $\hat v=v-x$ and call $v$ \emph{interior} if
$\hat v\in[1,n-1]^4$; let $I_x$ be the set of interior vertices. For $z\in I_x$
put $\tau(z)=(z-e_{\sigma(z)},\sigma(z))$, a link of $F$ with head $z$, and
$T_x=\{\tau(z):z\in I_x\}$. Every link with an end point in $I_x$ has its base
point in $Q_x$, so $T_x\subset V_x$, the map $\tau$ is injective, and no link
outside $V_x$ touches $I_x$. Call the links of $V_x\setminus T_x$ \emph{free}.
For $\psi\in G^{V_x\setminus T_x}$ let $[\psi,U]$ be the configuration equal to
$\psi$ on the free links, to $1$ on $T_x$ and to $U$ outside $V_x$, and let
$\gamma^T_{x,U}$ be the probability measure on $G^{V_x\setminus T_x}$ with
density proportional to $\prod_{p\cap V_x\ne\emptyset}\varphi([\psi,U]_p)$.

\begin{lemma}\label{lem:cubegauge}
For every bounded measurable $f\in\mathcal A$,
$\mathcal P_xf(U)=\int f([\psi,U])\,\gamma^T_{x,U}(d\psi)$.
\end{lemma}

\begin{proof}
Let $\sigma'\in G^{V_x}$ and write $(\sigma',U)$ for the configuration equal to
$\sigma'$ on $V_x$ and to $U$ elsewhere. Order $I_x$ by the sum of the
coordinates of $\hat z$; then $z-e_{\sigma(z)}$ precedes $z$ or is not
interior. Put $g_v=1$ for $v\notin I_x$ and recursively $g_z=g_{z-e_{\sigma(z)}}
\sigma'_{\tau(z)}$. The gauge transformation $g$ sets every link of $T_x$ to $1$
and fixes every link outside $V_x$, so $g\cdot(\sigma',U)=[\psi,U]$ with
$\psi=(g\cdot\sigma')|_{V_x\setminus T_x}$. Since $g$ depends only on
$\sigma'_{T_x}$, for fixed $\sigma'_{T_x}$ the map $\sigma'_{V_x\setminus T_x}
\mapsto\psi$ multiplies each coordinate on both sides by fixed elements, so
$\sigma'\mapsto(\psi,\sigma'_{T_x})$ preserves Haar measure. Plaquette variables
change by conjugation and $\varphi$ is a class function, so the density of
$\gamma^0_{V_x}(\cdot\mid U)$ at $\sigma'$ is proportional to
$\prod_{p\cap V_x\ne\emptyset}\varphi([\psi,U]_p)$, while $f(\sigma',U)=
f([\psi,U])$. Integrating out $\sigma'_{T_x}$ gives the claim.
\end{proof}

A link $(b,\mu)$ of $V_x$ is a \emph{face link} if $\hat b_i=0$ for some
$i\ne\mu$, that is, if both its end points lie in one lower face
$\{\hat v_i=0\}$ of the cube.

\begin{lemma}\label{lem:fifteen}
\begin{enumerate}[label=\textup{(\roman*)},leftmargin=*]
\item If all four links of a plaquette are free, its link of $F$ is a face
link.
\item Every free link has at most fifteen free partners.
\end{enumerate}
\end{lemma}

\begin{proof}
(i) Let $\phi=(y,s)$ be the link of $F$ of the plaquette $p$ and $h=y+e_s$.
If $h$ were interior, then $\sigma(h)=s$ and $\phi=\tau(h)\in T_x$. Since $y\in
Q_x$, $\hat h_i=\hat y_i\in[0,n-1]$ for $i\ne s$ and $\hat h_s\in[1,n]$, so
$\hat h_s=n$ or $\hat y_i=0$ for some $i\ne s$. If $\hat h_s=n$, the link of
$p$ parallel to its second direction $\nu$ through $h$ has a base point with
$s$-coordinate $n$ and is not in $V_x$. Hence $\phi$ is a face link.

(ii) Let $k=(b,\mu)$ be free and $m(p)$ the number of free links of a
plaquette $p\ni k$; $k$ has $\sum_{p\ni k}(m(p)-1)$ free partners. If $k$ is a
face link, $\hat b_i=0$ with $i\ne\mu$, the plaquette of $k$ with base point
$b-e_i$ in the plane $\{\mu,i\}$ has its other three links based at
$i$-coordinate $-1$, outside $V_x$; the other five plaquettes give at most $3$
each, $15$ in all. Otherwise $k$ has at most $12+M$ free partners, where $M$
counts the plaquettes $p\ni k$ with $m(p)=4$, whose link $\phi$ of $F$ is a face
link by (i), so $\phi\ne k$. If $k\in F$ then $M=0$; let $k$ be active and
$\phi$ lie in the face $\{\hat v_i=0\}$. If $\phi$ shares an end point $v$ with
$k$, then $v=b$ (for $v=b+e_\mu$ either $k$ is a face link or $\hat b_i=-1$),
$\hat b_i=0$ and $i=\mu$; there are two links of $F$ at $b$, each in at most
one plaquette of $k$. If $\phi$ is opposite to $k$, it is $(b\pm e_\nu,\mu)$
in a plaquette of the plane $\{\mu,\nu\}$; as $k$ is not a face link, $\nu=i$,
the sign is $-$ and $\hat b_i=1$, and by Lemma~\ref{lem:structure}(ii) the
plane $\{\mu,\nu\}$ is the one plane in which the link of $F$ is opposite to
$k$, so this happens for at most one plaquette. Hence $M\le3$.
\end{proof}

For a free link $e$ the conditional law $\gamma^T_e(\cdot\mid\psi,U)$ of
$\psi_e$ under $\gamma^T_{x,U}$ has density proportional to $\prod_{p\ni e}
\varphi([\psi,U]_p)$. Let $C^T_{ek}$ be the supremum of $W_1(\gamma^T_e(\cdot
\mid\psi,U),\gamma^T_e(\cdot\mid\psi',U'))/\delta([\psi,U]_k,[\psi',U']_k)$ over
pairs for which $[\psi,U]$ and $[\psi',U']$ differ only at a link $k$, free or
outside $V_x$; it vanishes unless $e$ and $k$ are partners. For $j\notin V_x$
let $b(V_x,j)$ be the number of links of $V_x$ that are partners of $j$, and
\[
R_x(j)=\sup\frac{W_1^{d_x}(\gamma^T_{x,U},\gamma^T_{x,U'})}{\delta(U_j,U'_j)},
\qquad d_x(\psi,\psi')=\sum_{e\in V_x\setminus T_x}\delta(\psi_e,\psi'_e),
\]
the supremum over configurations that differ only at $j$.

\begin{lemma}\label{lem:bpblock}
If $C^T_{ek}\le c$ for all free $e$ and all $k$, with $15c<1$, then
$R_x(j)\le c\,b(V_x,j)/(1-15c)$.
\end{lemma}

\begin{proof}
Let $U,U'$ differ only at $j$, let $\pi$ minimise $\int d_x\,d\pi$ over
couplings of $\gamma^T_{x,U}$ and $\gamma^T_{x,U'}$, and let $\varrho(e)=\int
\delta(\psi_e,\psi'_e)\,d\pi$. Replacing $(\psi_e,\psi'_e)$ by a measurably
chosen optimal coupling of the two conditional laws of $\psi_e$
\cite[Cor.~5.22]{Vil09} gives another coupling, so by minimality and the
triangle inequality $\varrho(e)\le c\sum_k\varrho(k)+c[e\text{ a partner of }j]
\delta(U_j,U'_j)$, the sum over the free partners $k$ of $e$. Summing over $e$
and using Lemma~\ref{lem:fifteen}(ii), $(1-15c)\sum_e\varrho(e)\le
c\,b(V_x,j)\delta(U_j,U'_j)$.
\end{proof}

\begin{lemma}\label{lem:bpcount}
For $n\ge2$, $L\ge2n+4$ and a link $j$, $\sum_{x:\,j\notin V_x}b(V_x,j)=
18n^3-3n^2$.
\end{lemma}

\begin{proof}
Let $b$ be the base point of $j$ and $\nu$ its direction. For each $\mu\ne\nu$
the partners in the two plaquettes of the plane $\{\mu,\nu\}$ have base points
$b+\Delta$, $\Delta\in\{e_\nu,e_\mu,0,-e_\mu,-e_\mu,e_\nu-e_\mu\}$. With $z=b-x$
and $Q=[0,n)^4$, a link based at $b+\Delta$ lies in $V_x$ exactly when $z+\Delta\in
Q$, and $j\in V_x$ when $z\in Q$. The sum is $\sum_\Delta(n^4-|Q\cap(Q-\Delta)|)$,
which is $0$ for $\Delta=0$, $n^3$ for $\Delta=\pm e_\mu,e_\nu$ and $2n^3-n^2$ for
$\Delta=e_\nu-e_\mu$, that is $6n^3-n^2$ for each of the three $\mu$.
\end{proof}

\begin{lemma}\label{lem:onelink}
\begin{enumerate}[label=\textup{(\roman*)},leftmargin=*]
\item For $G=\SU(2)$, the chordal metric and $\varphi=e^{2\beta\ReTr(\cdot)}$,
$C^T_{ek}\le|\beta|=|\beta_W|/4$.
\item Let $G$ have semisimple Lie algebra, $\delta$ the geodesic distance of
$\ip XY=-\Tr(XY)$, and $6N|\beta|<\kappa_G$. Then $C^T_{ek}\le N|\beta|/
(\kappa_G-6N|\beta|)$.
\end{enumerate}
\end{lemma}

\begin{proof}
(i) The conditional law of $\psi_e$ is the von Mises--Fisher law $\nu_b$ with
$b=4\beta\sum_{p\ni e}\bar s_p$, where $s_p$ is the product of the other three
variables of $p$, as in Sec.~\ref{sec:state}. A partner $k$ enters one
$s_p$, and moving it by $\delta_k$ moves $\bar s_p$ by $\delta_k$, since left and
right multiplication and inversion are isometries; the sharp quarter of
\cite{BB26a} gives the bound. Links of $T_x$ do not move.

(ii) The conditional law is $\nu\propto e^{-V}d\psi_e$ with $V=-N\beta
\sum_{p\ni e}\ReTr[\psi,U]_p$; by~\eqref{eq:fp} each term has second
derivative $\ReTr(Y^2W)$ with $|Y|=|X|$, so $\Ric_G+\Hess V\ge K_1=
\kappa_G-6N|\beta|$. Lemma~\ref{lem:weighted} with one factor and the argument
of Lemma~\ref{lem:decayrate} give $|\Cov_\nu(g,h)|\le\|\nabla g\|_\infty
\|\nabla h\|_\infty/K_1$. Move the partner $k$ along a unit-speed geodesic; the
derivative of $\nu_s(g)$ is $-\Cov_{\nu_s}(g,\partial_sV_s)$, and $\partial_sV_s$
comes from the one plaquette containing $e$ and $k$, with gradient in $\psi_e$
of norm at most $N|\beta|$ by~\eqref{eq:mixed}. For smooth $1$-Lipschitz $g$
this gives $|\nu_0(g)-\nu_\delta(g)|\le N|\beta|\delta/K_1$; convolution on $G$
preserves Lipschitz constants for a bi-invariant metric, and
Kantorovich--Rubinstein duality concludes.
\end{proof}

\begin{proposition}\label{prop:interiorbp}
If $\sum_{x:\,i\notin V_x}R_x(i)\le\alpha n^4$ for every link $i$, with
$\alpha<1$, then $\BP(n,1-\alpha)$ holds.
\end{proposition}

\begin{proof}
For a function $h$ let $\mathrm{Lip}(h)$ be the least constant with
$|h(U)-h(U')|\le\mathrm{Lip}(h)\delta(U_i,U'_i)$ whenever $U,U'$ differ only at
one link $i$; then $|h(U)-h(U')|\le\mathrm{Lip}(h)d_E(U,U')$. Let $h\in\mathcal A$
be Lipschitz and $U,U'$ differ only at $i$. If $i\in V_x$,
$\mathcal P_xh(U)=\mathcal P_xh(U')$. If $i\notin V_x$, Lemma~\ref{lem:cubegauge}
and an optimal coupling $\pi$ of $\gamma^T_{x,U}$ and $\gamma^T_{x,U'}$ give
\[
|\mathcal P_xh(U)-\mathcal P_xh(U')|\le\int|h([\psi,U])-h([\psi',U'])|\,d\pi
\le\mathrm{Lip}(h)\bigl(R_x(i)+1\bigr)\delta(U_i,U'_i).
\]
Since $i$ lies in $n^4$ of the cubes, $\mathrm{Lip}(\mathcal Kh)\le(1-\kappa_0)
\mathrm{Lip}(h)$ with $\kappa_0=(1-\alpha)n^4/L^4$, and as $\mathcal K$
preserves $\mathcal A$, $\Var_\mu(\mathcal K^th)\le\frac12(1-\kappa_0)^{2t}
\mathrm{Lip}(h)^2D^2|E|^2$. The restriction of $\mathcal K$ to the mean-zero
part of $\mathcal A$ is self-adjoint with $0\le\mathcal K\le I$, and
$\Var_\mu(\mathcal K^th)=\int\lambda^{2t}d\nu_h$ for its spectral measure $\nu_h$
at $h-\mu(h)$, so $\nu_h$ lives on $[0,1-\kappa_0]$. Lipschitz functions are
dense in $L^2(\mu)$, and averaging over the compact gauge group, the
orthogonal projection onto $\mathcal A$, does not increase $\mathrm{Lip}$,
because gauge transformations act link by link by isometries; so Lipschitz
functions are dense in $\mathcal A$, and $\ip f{(I-\mathcal K)f}\ge\kappa_0
\Var_\mu(f)$ for real $f\in\mathcal A$, which is $\BP(n,1-\alpha)$.
\end{proof}

\begin{theorem}\label{thm:bpstrong}
Let $n\ge2$ and $L\ge2n+4$. If $C^T_{ek}\le c$ for all cubes, free links $e$
and links $k$, with $15c<1$, then $\BP(n,1-\alpha_n)$ holds whenever
\[
\alpha_n=\frac{c\,(18n-3)}{(1-15c)\,n^2}<1 .
\]
In particular:
\begin{enumerate}[label=\textup{(\roman*)},leftmargin=*]
\item for $G=\SU(2)$ and $|\beta_W|<\frac4{15}$, with $\alpha_n=|\beta_W|(18n-3)/
((4-15|\beta_W|)n^2)$;
\item for $G\subset\U(N)$ with semisimple Lie algebra and $21N|\beta|<\kappa_G$,
with $\alpha_n=N|\beta|(18n-3)/((\kappa_G-21N|\beta|)n^2)$.
\end{enumerate}
For fixed $\beta$ in these ranges $\alpha_n$ decreases to $0$, so $\BP(n,\varpi)$
holds with $\varpi\to1$ as $n\to\infty$, uniformly in the torus.
\end{theorem}

\begin{proof}
Lemmas~\ref{lem:bpblock} and~\ref{lem:bpcount} give $\sum_{x:\,i\notin V_x}
R_x(i)\le c(18n^3-3n^2)/(1-15c)=\alpha_nn^4$, and
Proposition~\ref{prop:interiorbp} applies. Lemma~\ref{lem:onelink} gives
$c=|\beta_W|/4$ for $\SU(2)$, and $c=N|\beta|/(\kappa_G-6N|\beta|)$ in
general, for which $15c<1$ reads $21N|\beta|<\kappa_G$ and $c/(1-15c)=
N|\beta|/(\kappa_G-21N|\beta|)$.
\end{proof}

For $\SU(2)$ the condition $\alpha_n<1$ reads $|\beta_W|(15n^2+18n-3)<4n^2$; at
$|\beta_W|=\frac14$ it holds for $n\ge18$, and its threshold increases to
$\frac4{15}$ as $n\to\infty$. At strong coupling the condition is therefore
\emph{easier} for large cubes, because the response of a cube to one boundary
link is spread over $n^4$ links while it touches only $O(n^3)$ of them.

\subsection*{A weighted count}
For a free link $e$ let $N(e)$ be its set of free partners. The proof of
Lemma~\ref{lem:bpblock} gives, for each free $e$,
\begin{equation}\label{eq:rho}
\varrho(e)\le\sum_{k\in N(e)}C^T_{ek}\varrho(k)+C^T_{ej}[e\text{ a partner of }j]
\,\delta(U_j,U'_j).
\end{equation}

\begin{lemma}\label{lem:weightedresp}
Let $w$ be a function on the free links with values in $[w_-,w_+]$, $w_->0$,
with $\sum_{k\in N(e)}w(k)\le\lambda w(e)$ for every free $e$. If $C^T_{ek}\le
c$ for all free $e$ and all $k$, and $c\lambda<1$, then
$R_x(j)\le(w_+/w_-)\,c\,b(V_x,j)/(1-c\lambda)$.
\end{lemma}

\begin{proof}
Multiply~\eqref{eq:rho} by $w(e)$ and sum over $e$. Since $k\in N(e)$ exactly
when $e\in N(k)$, $\sum_ew(e)\varrho(e)\le c\lambda\sum_kw(k)\varrho(k)+
c\,w_+b(V_x,j)\delta(U_j,U'_j)$, and $\sum_e\varrho(e)\le w_-^{-1}\sum_ew(e)
\varrho(e)$.
\end{proof}

Write a link of $V_x$ as $(x+b,\mu)$ with $b\in[0,n-1]^4$, let
$z(e)=\#\{i\ne\mu:b_i=0\}$ be the number of lower faces containing it, and
call it a \emph{foot} if $b_\mu=0$. On free links put $w=45$ if $z=1$ and $e$
is active and not a foot, $w=47$ if $z=1$ and $e\in F$ or $e$ is a foot,
$w=48$ if $z=0$ and $e$ is an active foot, and $w=40$ otherwise.

\begin{proposition}\label{prop:weights}
For every $n\ge2$, every cube and every free link $e$,
$\sum_{k\in N(e)}w(k)\le\frac{72}5w(e)$.
\end{proposition}

\begin{proof}
Work in the coordinates $b$, so that $(y,\nu)\in F$ when $\sigma(x+y)=\nu$.
Let $e=(b,\mu)$. For $\nu\ne\mu$ the other links of the plaquette $p^+_\nu$ of
base point $b$ in the plane $\{\mu,\nu\}$ are $P_1=(b+e_\mu,\nu)$,
$P_2=(b+e_\nu,\mu)$, $P_3=(b,\nu)$, and those of the plaquette $p^-_\nu$ of base
point $b-e_\nu$ are $P_4=(b-e_\nu,\mu)$, $P_5=(b-e_\nu+e_\mu,\nu)$,
$P_6=(b-e_\nu,\nu)$. If $e$ is active, with $s=\sigma(x+b)$,
$t=\sigma(x+b+e_\mu)$ and $\rho$ the fourth direction, the links of $F$ among
them are $P_3,P_6$ for $\nu=s$, $P_1,P_5$ for $\nu=t$ and $P_2,P_4$ for
$\nu=\rho$: the first four by Lemma~\ref{lem:structure}(i) at $b$ and at
$b+e_\mu$, and in the plane $\{\mu,\rho\}$ the link of $F$ is opposite to $e$
by Lemma~\ref{lem:structure}(ii). If $e\in F$, none of them is in $F$.

\emph{Reduction to the orthant.} Call a link with base point in
$O=\Z^4_{\ge0}$ \emph{$O$-free} if it is active or has $z\ge1$, extend $w$ to
such links by the same rule, and count as partners of $e$ the $O$-free links
of the plaquettes of $e$ with base point in $O$. If $b_\nu=0$, the links of
$p^-_\nu$ other than $e$ lie outside $V_x$. A link $(y,\mu')\in V_x$ lies in
$T_x$ exactly when it is in $F$ and its head is interior, that is, $z=0$ and
$y_{\mu'}\le n-2$; so free links are $O$-free except the \emph{exceptional}
ones, links of $F$ with $z=0$ and $y_{\mu'}=n-1$, of weight $40$. An
exceptional link $\phi$ in a plaquette $p$ is its link of $F$, and $p$ has a
link $\ell$ based at $\mu'$-coordinate $n$, which lies outside $V_x$, is
active, has its base point in $O$ (as $z(\phi)=0$), and has $w(\ell)\ge40$;
exchanging $\phi$ for $\ell$ shows that the orthant sum dominates the sum over
$N(e)$ for non-exceptional $e$. An exceptional $e$ has at most two partners in
$V_x$ in each plaquette, each of weight at most $48$, and $12\cdot48=576=
\frac{72}5\cdot40$.

\emph{The orthant count.} Fix $\nu\ne\mu$, put $f=[b_\mu=0]$, $\beta=b_\nu$,
and let $z'$ be the number of zeros among $b_i$, $i\notin\{\mu,\nu\}$. From
the base points, $z(P_1)=z(P_2)=z(P_5)=z'$, $z(P_3)=z(P_6)=z'+f$,
$z(P_4)=z'+[\beta=1]$; $P_2,P_4$ are feet when $f=1$, $P_1,P_3$ when
$\beta=0$, $P_5,P_6$ when $\beta=1$; and $p^-_\nu$ counts only when
$\beta\ge1$. An $O$-free link has weight $40$ if $z\ge2$, $47$ or $45$ if
$z=1$ according as it is in $F$ or a foot or neither, and $48$ or $40$ if
$z=0$ (then it is active) according as it is a foot or not; a link of $F$ with
$z=0$ contributes nothing. Let $C(\text{role},f,\beta,z')$ be the resulting
weight of the $O$-free links among $P_1,\dots,P_6$, the role of $\nu$ being
$s$, $t$, $\rho$, or $\mathrm F$ when $e\in F$. For $z'=2$, $C=120$ if
$\beta=0$ and $C=240$ otherwise; for $z'\le1$ the values are
\begin{center}
\small
\begin{tabular}{cc|rrr|rrr}
\toprule
&&\multicolumn{3}{c|}{$z'=0$, $\beta=$}&\multicolumn{3}{c}{$z'=1$, $\beta=$}\\
role&$f$&$0$&$1$&${\ge}2$&$0$&$1$&${\ge}2$\\
\midrule
$s$&0&88&173&160&139&271&274\\
$s$&1&143&277&270&134&259&264\\
$t$&0&88&173&160&139&271&274\\
$t$&1&95&187&186&134&261&268\\
$\rho$&0&96&223&160&141&271&274\\
$\rho$&1&95&227&170&134&259&264\\
$\mathrm F$&0&136&261&240&139&269&270\\
$\mathrm F$&1&143&275&266&134&259&264\\
\bottomrule
\end{tabular}
\end{center}
For instance, for the role $s$, $f=1$, $\beta=1$, $z'=0$ the weights of
$P_1,\dots,P_6$ are $40,48,47,47,48,47$. The orthant sum is
$\sum_{\nu\ne\mu}C(\text{role}(\nu),f,b_\nu,z'_\nu)$, where $z'_\nu=z(e)-1$ if
$b_\nu=0$ and $z'_\nu=z(e)$ otherwise, and $\frac{72}5w(e)$ is $576$, $648$,
$676.8$ or $691.2$ for $w(e)=40,45,47,48$. For active $e$ with $z=0$ the sum
is at most $173+173+223=569$ ($f=0$, $w=40$) or $277+187+227=691$ ($f=1$,
$w=48$). For active $e$ with $z=1$ it is at most $96+274+274=644$ ($f=0$,
$w=45$), and for $f=1$, $w=47$, at most $143+268+264$, $95+264+264$ or
$95+264+268$ as the zero direction has role $s$, $t$ or $\rho$, all at most
$675$. For active $e$ with $z=2$ it is at most $141+139+240=520$, for $z=3$
it is $360$, with $w=40$. For $e\in F$ (then $z\ge1$) it is at most
$136+270+270=676$ or $143+264+264=671$ when $z=1$, $w=47$; at most $518$ when
$z=2$; and $360$ when $z=3$, with $w=40$.
\end{proof}

\begin{theorem}\label{thm:bpweighted}
Let $n\ge2$ and $L\ge2n+4$. If $C^T_{ek}\le c$ for all cubes, free links $e$
and links $k$, with $\frac{72}5c<1$, then $\BP(n,1-\alpha_n)$ holds whenever
$\alpha_n=6c(18n-3)/((5-72c)n^2)<1$. In particular:
\begin{enumerate}[label=\textup{(\roman*)},leftmargin=*]
\item for $G=\SU(2)$ and $|\beta_W|<\frac5{18}$, with
$\alpha_n=3|\beta_W|(18n-3)/(2(5-18|\beta_W|)n^2)$;
\item for $G=\SU(3)$, the Hilbert--Schmidt distance and $0\le\beta_W\le0.45$.
\end{enumerate}
\end{theorem}

\begin{proof}
Lemma~\ref{lem:weightedresp} with Proposition~\ref{prop:weights},
$\lambda=\frac{72}5$ and $w_+/w_-=\frac65$, and Lemma~\ref{lem:bpcount} give
$\sum_{x:\,i\notin V_x}R_x(i)\le\frac65c(18n^3-3n^2)/(1-\frac{72}5c)=
\alpha_nn^4$, and Proposition~\ref{prop:interiorbp} applies; its proof uses
only that $\delta$ is bi-invariant and invariant under inversion. For (i),
$c=|\beta_W|/4$ by Lemma~\ref{lem:onelink}. For (ii) let $r=2\beta_W$,
$Z_*(r)=\int_{\SU(3)}e^{r\ReTr U}dU$, $\epsilon(r)=1-5e^{-r/5}/(9Z_*(r))$ and
$\gamma(r)=\frac{283}{128}-\frac r2$. The conditional law of a free link is
$\nu_b\propto e^{\ReTr(Ub)}dU$ with $b$ a sum of $N\beta=r/6$ times six
products of elements of $\SU(3)$, the links of $T_x$ being $1$, so the
one-link estimate of the $\SU(3)$ theorem of~\cite{BBLRS26} applies and gives
$c=r\sqrt{\epsilon(r)}/(6\gamma(r))$, increasing in $r$. With the moments
$m_k=\int(\ReTr U)^kdU$, $(m_0,\dots,m_8)=(1,0,\frac12,\frac14,\frac34,
\frac{15}{16},\frac{65}{32},\frac{231}{64},\frac{973}{128})$, and $|m_k|\le
3^{k-8}m_8$ for $k\ge8$,
\[
Z_*(r)\le\sum_{k=0}^8\frac{m_kr^k}{k!}+\frac{m_8}{3^8}\frac{(3r)^9}{9!}
\Bigl(1-\frac{3r}{10}\Bigr)^{-1}.
\]
At $r=0.9$ the sum is at most $1.2599146$ and the tail at most
$3.4\cdot10^{-5}$, so $Z_*(0.9)\le1.2599486$; with $e^{-0.18}\ge1-0.18+0.0162
-0.000972$ this gives $\epsilon(0.9)\le0.63172<0.7949^2$, and with
$\gamma(0.9)=\frac{1127}{640}$, $\frac{72}5c\le2.16\cdot0.7949\cdot640/1127
\le0.9751$.
\end{proof}

For fixed $n$ the threshold of (i) exceeds that of
Theorem~\ref{thm:bpstrong}(i) exactly when $n^2-6n+1>0$, that is for $n\ge6$;
both theorems hold. The same arithmetic at $r=0.88$, where $Z_*\le1.246559$,
$\epsilon\le0.62627$ and $\gamma=\frac{5667}{3200}$, gives $15c\le0.9832$, so
Theorem~\ref{thm:bpstrong} holds for $\SU(3)$ with $0\le\beta_W\le0.44$.

Without the gauge fixing every link has eighteen partners, and the same
argument needs $18c<1$, that is $|\beta_W|<\frac29$ for $\SU(2)$. No choice of
gauge-fixed links and weights goes much further:

\begin{proposition}\label{prop:methodlimit}
Let $T\subset V_x$ have at most $(n-1)^4$ links and $A$ be the partner matrix
of $V_x\setminus T$. Then $\lambda_{\max}(A)\ge6r(4-r)(3-r)/(4-r^2)$ with
$r=(1-1/n)^2$, which tends to $12$. If $c\lambda_{\max}(A)\ge1$, the
inequalities $0\le\varrho\le cA\varrho+s$ used in Lemma~\ref{lem:bpblock}
admit solutions with $\sum_e\varrho(e)$ arbitrarily large, so for $\SU(2)$ the
argument cannot pass $|\beta_W|=\frac13$.
\end{proposition}

\begin{proof}
The $P=6n^2(n-1)^2=6rn^4$ plaquettes with all links in $V_x$ contain at most
$6|T|=aP$ links of $T$ in all, $a\le r$, so their numbers $m(p)$ of links outside
$T$ have mean at least $4-a\ge3$; by convexity $\mathbf1^{\top}A\mathbf1\ge
\sum_pm(p)(m(p)-1)\ge P(4-a)(3-a)$, and the Rayleigh quotient of $\mathbf1$ is
at least $6r(4-a)(3-a)/(4-ar)$, which decreases in $a$ because
$r/(4-ar)\le1/(4-a)$. For a Perron vector $v\ge0$ of $A$, $\varrho=tv$
satisfies $cA\varrho\ge\varrho$ for every $t>0$.
\end{proof}

\section{Free massive fields}\label{sec:free}

Here $d\ge1$, $n\ge2$, $L\ge n+1$, $\mathbb T=(\Z/L\Z)^d$, $\mu>0$ is a lattice
mass, $-\Delta$ is the lattice Laplacian and $Q=-\Delta+\mu^2$. For a positive
definite $Q$ on $\R^N$ let $\gamma_Q$ be the centred Gaussian law with
covariance $C=Q^{-1}$, and for $V\subset\{1,\dots,N\}$ let $\mathcal P_V$ be the
conditional expectation given the coordinates outside $V$. For the cubes
$V_x=x+\{0,\dots,n-1\}^d$, $x\in\mathbb T$, put
\begin{equation}\label{eq:ffvarpi}
\varpi(Q)=\inf_f\frac{n^{-d}\sum_x\|f-\mathcal P_{V_x}f\|^2}{\Var_{\gamma_Q}(f)},
\end{equation}
the infimum over nonconstant $f\in L^2(\gamma_Q)$, the best constant in
$\BP(n,\cdot)$. With $s=n\mu$ and $\mu=ma$, $n=\ell/a$, the parameter $s=m\ell$
is the mass times the physical size of the cube.

\begin{lemma}[linear functions suffice]\label{lem:chaos}
For a positive definite $Q$ and any family of coordinate sets $(V_x)$, the
infimum in~\eqref{eq:ffvarpi} equals its infimum over the linear functions
$\ell_h(\phi)=\ip h\phi$, and
\[
\|\ell_h-\mathcal P_V\ell_h\|^2=\ip{h_V}{Q_{VV}^{-1}h_V},\qquad
\Var(\ell_h)=\ip h{Q^{-1}h},
\]
with $h_V$ the restriction of $h$ to $V$ and $Q_{VV}$ the block of $Q$ on $V$.
\end{lemma}

\begin{proof}
The conditional law of $\phi_V$ given $\phi_{V^c}$ is Gaussian with covariance
$Q_{VV}^{-1}$, which gives the formula for $\ell_h$. Write $\phi=Q^{-1/2}\xi$ with
$\xi$ standard Gaussian. Conditioning on $\phi_{V^c}$ is conditioning on the
Gaussian subspace $K_V=\{\ip u\xi:u\in Q^{-1/2}\R^{V^c}\}$, and on the $m$-th
Wiener chaos, identified with the symmetric tensors $\mathrm{Sym}^m(\R^N)$, it
acts as $\pi_V^{\otimes m}$, with $\pi_V$ the orthogonal projection onto $Q^{-1/2}
\R^{V^c}$~\cite[Ch.~4]{Jan97}. The chaoses are orthogonal and invariant under
every $\mathcal P_V$, so it suffices to treat one chaos $m\ge1$. The projection
$I-\pi_V^{\otimes m}$ dominates $(I-\pi_V)\otimes I^{\otimes(m-1)}$, as both are
projections that commute and the range of the second lies in that of the
first. For a symmetric tensor $F$, therefore,
\[
\sum_x\|F-\pi_{V_x}^{\otimes m}F\|^2\ge\Bigl\langle F,\Bigl(\sum_x(I-\pi_{V_x})\Bigr)
\otimes I^{\otimes(m-1)}F\Bigr\rangle\ge\lambda\|F\|^2,
\]
where $\lambda$ is the least eigenvalue of $\sum_x(I-\pi_{V_x})$ on $\R^N$, which is
the infimum of the quotient over linear functions ($m=1$).
\end{proof}

\begin{theorem}[the exact constant]\label{thm:free}
Let $V=\{0,\dots,n-1\}^d$, $-\Delta_D$ its Dirichlet Laplacian and $\mathbf 1$ its
constant function. Then
\[
\varpi(-\Delta+\mu^2)=\frac{\mu^2}{n^d}\ip{\mathbf 1}{(-\Delta_D+\mu^2)^{-1}\mathbf 1},
\]
attained only at multiples of $\ell_{\mathbf 1}$ among linear functions, and
\[
\Bigl(\frac56\Bigr)^d\frac{s^2}{s^2+10d}\le\varpi(-\Delta+\mu^2)\le\frac{s^2}
{s^2+4dn^2\sin^2\frac\pi{2n+2}}\le\frac{s^2}{s^2+16d/9}.
\]
\end{theorem}

\begin{proof}
Since $L\ge n+1$, the block $Q_{V_xV_x}$ is the translate of $Q_{VV}=-\Delta_D+
\mu^2$; put $G_V=Q_{VV}^{-1}$. By Lemma~\ref{lem:chaos}, $n^d\varpi$ is the least
value of $\ip h{Sh}/\ip h{Q^{-1}h}$ with $S=\sum_xR_{V_x}^{\mathsf T}Q_{V_xV_x}^{-1}
R_{V_x}$, that is, the least eigenvalue of $Q^{1/2}SQ^{1/2}$. Translations commute
with $Q$ and $S$, so the plane waves $e_k(y)=e^{ik\cdot y}$ diagonalise it with
eigenvalues $\lambda_k=\hat q(k)\ip{e_k^V}{G_Ve_k^V}$, where $\hat q(k)=\mu^2+
\sum_i(2-2\cos k_i)$ and $e_k^V$ is the restriction to $V$.

Let $B$ map $f$ to the function $u\mapsto\sum_{w\sim u,\,w\notin V}f(w)$ on $V$.
Then $Q_{VV}R_Vf=R_VQf+Bf$, so $Q_{VV}e_k^V=\hat q(k)e_k^V+Be_k$, and pairing
$G_V$ of this with $e_k^V$ gives
\[
\lambda_k=n^d-\sum_{u,u'\in V}\sum_{w\sim u',\,w\notin V}G_V(u,u')\cos\bigl(k\cdot
(w-u)\bigr).
\]
Write $Q_{VV}=(2d+\mu^2)(I-K)$ with $K$ entrywise nonnegative and $\|K\|\le
2d/(2d+\mu^2)<1$; the Neumann series shows that $G_V$ has positive entries.
Hence $\lambda_k-\lambda_0=\sum G_V(u,u')(1-\cos(k\cdot(w-u)))\ge0$, with equality only
at $k=0$ (take $u=u'$ on a face and $w=u+e_i$). Since $Q\mathbf 1=\mu^2\mathbf 1$,
$\lambda_0=\mu^2\ip{\mathbf 1}{G_V\mathbf 1}$, which is the formula.

\emph{Upper bound.} $G_V\le(\lambda_1+\mu^2)^{-1}$ with $\lambda_1=4d\sin^2\frac\pi
{2n+2}$ the least eigenvalue of $-\Delta_D$, and $\|\mathbf 1\|^2=n^d$; then
$\sin t\ge2t/\pi$ gives $n^2\lambda_1\ge4dn^2/(n+1)^2\ge16d/9$.

\emph{Lower bound.} By the variational formula $\ip{\mathbf 1}{G_V\mathbf 1}\ge
\ip{\mathbf 1}w^2/\ip w{Q_{VV}w}$, take $w(u)=\prod_i\tau(u_i)$ with the parabolic
profile $\tau(s)=6(s+1)(n-s)/(n(n+1)(n+2))$ on $\{0,\dots,n-1\}$, zero
outside. Then $\sum\tau=1$, $S_2=\sum\tau^2=\alpha_n/n$ and $D_2=\sum_{s\in\Z}
(\tau(s+1)-\tau(s))^2=\beta_n/n^3$ with
\[
\alpha_n=\frac{6(n^2+2n+2)}{5(n+1)(n+2)}\le\frac65,\qquad
\beta_n=\frac{12n^2}{(n+1)(n+2)}=10\rho_n\alpha_n,\quad
\rho_n=\frac{n^2}{n^2+2n+2}<1,
\]
by the power sums $\sum_{u=1}^nu(n+1-u)=n(n+1)(n+2)/6$, $\sum_{u=1}^nu^2(n+1-u)^2=
n(n+1)(n+2)(n^2+2n+2)/30$ and $\sum_{u=0}^n(n-2u)^2=n(n+1)(n+2)/3$. The Dirichlet
form of the product is $\ip w{-\Delta_Dw}=dS_2^{d-1}D_2$ and $\ip{\mathbf 1}w=1$, so
\[
\varpi\ge\frac{\mu^2}{n^d(\mu^2S_2^d+dS_2^{d-1}D_2)}=\frac1{\alpha_n^d+d\alpha_n^{d-1}
\beta_ns^{-2}}=\alpha_n^{-d}\frac{s^2}{s^2+10d\rho_n}\ge\Bigl(\frac56\Bigr)^d
\frac{s^2}{s^2+10d}.\qedhere
\]
\end{proof}

At $n=2$ the two bounds coincide with $s^2/(s^2+4d)$. For $d=4$ the lower
bound is $0.0118$ at $s=1$ and $0.0438$ at $s=2$; the upper bound vanishes as
$s\to0$, so a massless free field satisfies no cube condition uniform in the
lattice spacing. The same proof, with the Hodge identity of the torus, gives
the cube condition for the free massive vector (Proca) field on gauge-invariant
functions, with the same lower bound and the upper bound $s^2/(1+s^2)$
\cite{BBLRS26}.
